%%
% BIThesis 本科毕业设计论文模板 —— 使用 XeLaTeX 编译 The BIThesis Template for Undergraduate Thesis
% This file has no copyright assigned and is placed in the Public Domain.
%%

\chapter{需求分析}

\section{总体需求分析}

本系统面向前端初学者构建移动客户端学习环境，目标是在移动设备上形成课程学习、代码练习、闯关挑战、每日任务、社交激励与个人成长相结合的完整学习闭环。根据需求文档，系统由移动客户端、服务端与管理后台端组成，其中移动客户端承担主要学习活动，服务端负责统一业务处理与数据支撑，管理后台端负责内容维护、运营支撑与数据管理。该类系统既需要满足移动学习场景下的便捷访问需求，也需要兼顾编程教育中实践反馈及时、任务组织明确和学习持续性较强等特点\cite{robins2003,heines2001,sommerville2015}。

从学习目标上看，系统不仅要支持前端基础知识的呈现和阅读，还应支持学习者通过编码练习和挑战任务完成知识内化，并借助成长反馈与社交激励增强持续学习意愿。因此，系统需求不能只停留在内容展示层面，而需要覆盖学习行为记录、任务组织、结果反馈和成长激励等多个方面\cite{hamari2014,ryan2000,kalmpourtzis2018}。前端基础知识包括 HTML 标记语言、CSS 样式语言和基础页面布局等内容，这些知识既具有明确的语法规则，又需要通过大量实践才能熟练掌握，因此系统需求特别强调理论与实践的结合。

从用户群体特征来看，目标用户主要是前端初学者，他们可能具备一定的计算机使用经验，但对 Web 开发技术尚不熟悉。这类用户在学习过程中容易因概念抽象、错误难以定位或缺乏及时反馈而产生挫败感，因此系统需要在交互设计上注重引导性、反馈的及时性和错误提示的可理解性。同时，初学者通常没有固定的学习时间和场所，移动学习场景正好契合其利用碎片时间学习的需求，但也对系统的离线能力和状态恢复能力提出了更高要求。

从系统边界来看，本文聚焦于学习应用本身的设计与实现，不包括完整的在线代码执行环境、大规模用户行为分析系统或商业化运营平台。AI 辅助功能限定在错误解释和提示层面，不扩展为完整的智能辅导系统。社交功能限定在好友、动态和排行榜等基础互动形式，不扩展为实时通信或协作编程等复杂场景。这些边界界定使系统需求保持在一个可实现的范围内，同时为核心学习闭环的完整性提供保障。

\section{技术需求分析}

在运行平台方面，移动客户端需要支持 iOS 与 Android 智能手机，适配不同尺寸屏幕并保持基础交互一致性；管理后台端主要面向桌面浏览器环境，用于内容维护与配置操作；服务端需要具备面向多端的统一 API 提供能力。由于系统存在移动客户端学习场景，因此网络环境不能假设始终稳定，应用需要具备一定的离线感知与本地缓存能力\cite{liang2018,jiao2023}。

移动客户端的技术需求还包括对移动设备特性的适配。例如，需要处理不同屏幕尺寸和分辨率的布局适配问题，需要在虚拟键盘弹出时调整界面元素位置，需要处理应用前后台切换时的状态保存和恢复。在性能方面，移动客户端需要控制内存占用和电量消耗，避免长时间使用导致设备过热或电量快速下降。在存储方面，移动客户端需要合理利用本地存储空间，对缓存数据进行有效管理，避免因存储膨胀影响应用性能。

在系统协同方面，服务端接口必须能够同时支撑移动客户端与管理后台端的访问需求，并通过统一响应结构保障调用一致性。与此同时，接口文档需要作为开发与联调的共同依据，以减少多端协作中的理解偏差。对移动客户端而言，还需要具备认证信息持久化、进度缓存、弱网回退与错误处理能力，以保证学习流程的连续性\cite{brusilovsky2007,papamitsiou2014,sommerville2015}。

服务端的技术需求还包括对并发访问的处理能力、对数据库操作的事务保障、对异常情况的容错处理等。虽然本文不涉及大规模性能优化，但服务端需要保证在正常负载下能够稳定响应请求，在异常情况下能够给出合理的错误信息而不是直接崩溃。管理后台端的技术需求则侧重于内容编辑的便利性、数据展示的清晰性和操作反馈的及时性，使运营人员能够高效地完成日常维护工作。

\section{功能需求分析}

结合需求文档与现有实现，移动客户端功能可归纳为以下几个方面。第一是账户管理功能，包括注册、登录与个人资料维护，用于建立用户身份与学习档案。第二是课程学习功能，包括课程列表浏览、课程详情查看、章节阅读和学习进度跟踪，用于支撑基础知识学习过程。第三是编码练习功能，包括练习详情查看、代码提交与结果反馈，用于连接知识理解与实践操作\cite{robins2003,heines2001,guo2013}。

第四是挑战相关功能，包括闯关挑战、星级评价、奖励发放与每日挑战，用于将学习活动组织为具有目标感与节奏感的任务链。第五是 AI 辅助功能，用于在学习者练习受阻时提供错误解释与必要提示，但其范围应限定在辅助解释层面，而不扩展为完整的代码生成或个性化推荐系统。第六是社交互动功能，包括好友、动态与排行榜，用于提供外部激励与学习氛围。第七是成长反馈功能，包括经验值、徽章、等级与统计信息，用于沉淀学习成果并强化用户成长感知\cite{koedinger1997,woolf2010,liang2018,jiao2023}。

在账户管理功能中，注册流程需要收集用户基础信息并验证输入有效性，登录流程需要验证身份并发放认证凭证，个人资料维护需要支持昵称、头像等信息的修改。这些功能虽然基础，但直接影响用户首次使用体验和后续身份识别，因此需要在易用性和安全性之间取得平衡。

在课程学习功能中，课程列表需要支持分类筛选和搜索，课程详情需要展示完整的课程信息和章节目录，章节阅读需要提供舒适的内容展示和进度记录。学习进度跟踪需要记录每门课程、每个章节的学习状态，使学习者能够随时了解自己的学习进展，也为系统的个性化推荐和成长反馈提供数据基础。

在编码练习功能中，题目展示需要清晰说明要求和示例，代码编辑需要提供基础的输入环境，结果反馈需要及时返回编译和测试结果。这些功能的实现质量直接影响学习者的练习体验，特别是错误反馈的可理解性，对于帮助初学者定位和修正错误至关重要。

对于支撑系统而言，管理后台端需具备课程、题目、挑战、用户与公告等内容管理能力，以支持移动客户端的内容供给与日常运营；服务端则需提供统一的认证、内容访问、提交记录、奖励结算、进度管理与社交数据接口，从而保证学习闭环能够在多端环境中稳定运行\cite{sommerville2015,campos2022}。

\section{非功能需求分析}

在可用性方面，系统应保证学习者能够以较短路径完成主要学习任务，界面反馈应清晰，状态切换应明确，减少因复杂操作造成的学习中断。具体而言，从打开应用到进入课程学习的操作步骤应控制在三步以内，常用功能应在主界面直接可达，避免深层嵌套的导航结构。错误提示应使用学习者能够理解的语言，避免展示技术性错误代码或堆栈信息\cite{heines2001,guo2013}。

在响应性方面，系统需要保证课程加载、页面切换、请求返回与错误反馈具备基本可接受的交互速度，并能在网络异常时给出可理解的提示。课程列表和章节内容的加载时间应控制在用户可接受的范围内，代码提交后的结果反馈应尽可能及时，避免学习者因等待时间过长而产生焦虑。在网络异常时，系统应明确告知用户当前网络状态，并提供重试或稍后操作的选项，而不是无限等待或显示空白页面。

在可维护性方面，移动客户端、服务端与管理后台端应保持职责清晰，接口结构统一，便于后续功能扩展与问题定位。代码组织应遵循模块化原则，各功能模块之间的耦合度应控制在合理范围内，使单个模块的修改不会对其他模块产生意外影响。文档应覆盖关键设计决策和接口定义，使新加入的开发者能够快速理解系统结构。

在一致性方面，系统需要通过统一契约文档控制接口字段、响应包络与权限边界，保证多端协作不因接口漂移而失效。移动客户端的界面风格应保持一致，相同功能的交互方式应在不同页面中遵循相同模式，减少学习者的认知负担。服务端对同类请求的处理逻辑应保持一致，避免因特殊处理导致的意外行为。

在可靠性方面，移动客户端应支持本地存储和离线回退策略，使学习记录不会因为短时断网而直接丢失。关键操作如代码提交应在本地保存副本，网络恢复后自动重试提交，避免因网络抖动导致学习成果丢失。服务端应保证数据持久化的可靠性，对关键操作采用事务保障，避免部分更新导致的数据不一致。

在扩展性方面，系统需要为后续课程扩充、挑战增加和能力增强保留一定结构空间，但本文只围绕当前实现范围展开，不对后续版本功能进行详细设计\cite{sommerville2015,papamitsiou2014}。具体而言，课程和挑战的数据模型应支持新增字段而不影响现有功能，接口设计应预留版本控制机制以便后续演进，功能模块的划分应使新增功能能够以插件方式接入而不需要大规模重构现有代码。

在兼容性方面，移动客户端需要适配主流 iOS 和 Android 系统版本，在保证核心功能可用的前提下，对较旧版本系统提供降级体验。服务端接口设计应遵循向后兼容原则，避免因接口变更导致旧版本移动客户端无法正常使用。管理后台端需要兼容主流桌面浏览器，确保运营人员能够在常见浏览器环境中顺利完成内容管理工作。

在可访问性方面，系统应遵循基本的可访问性设计原则，包括界面元素具有足够的对比度、交互控件具有适当的尺寸、状态变化具有明确的视觉反馈等。虽然本文不将可访问性作为核心研究目标，但这些基础考虑有助于扩大系统的适用人群，使更多学习者能够顺利使用系统进行前端学习。

在数据安全与隐私保护方面，系统需要妥善处理用户注册信息、学习行为数据和个人资料等敏感内容。用户密码在服务端以哈希形式存储，避免明文保存带来的泄露风险。学习行为数据仅供用户本人和授权管理员查看，不对外公开或用于未经授权的目的。在系统部署过程中，敏感信息应通过 HTTPS 等加密通道传输，降低数据在传输过程中被截获的风险。虽然本文不涉及专业的安全认证和隐私合规审计，但这些基础措施为系统的数据安全提供了必要的保障。

在国际化与本地化方面，当前系统主要面向中文用户群体，界面语言和内容均以中文呈现。结合本科毕业设计的实现边界，本文当前并未对多语言支持进行具体实现，而是将其作为后续可能的扩展方向保留讨论空间。
